\subsection{定理1的证明：第一部分}

蕴含关系 (i) \(\Rightarrow\) (ii) 已在（X，第 157 页，命题 5）中证明。蕴含关系 (ii) \(\Rightarrow\) (iii) 是显然的；对 (iv) \(\Rightarrow\) (v) 同理，因为 \(\beta_{\mathbf{M}}^{x}\) 与 \(\alpha_{\mathbf{M}}^{x} \otimes 1_{\mathbf{A}/\mathfrak{x}}\) 等同。

为证明 (iv) 蕴含 (iii)，考虑在 1 次分量上诱导的同态 \((\alpha_{M}^{x})_{1}\) 。设 E 表示分次 A-模 \(A[X_{1},\ldots,X_{n}]\) ；A-模 \(E_{1}\) 以 \(X_{1},\ldots,X_{n}\) 为基，而 \(\mathfrak{d}_{1}\) 是 \(E_{1}\) 的子 A-模，由元素 \((x_{i}X_{j}-x_{j}X_{i})\) （ \(1\leqslant i<j\leqslant n\) ）生成。因此 \(((E/\mathfrak{d})\otimes_{A}M)_{1}\) 与 \(K_{1}(x,M)/B_{1}(K.(x,M))\) 等同，同态 \((\alpha_{M}^{x})_{1}\) 等同于从 \(K_{1}(x,M)/B_{1}(K.(x,M))\) 到 \(B_{0}(K.(x,M))\) 上的、由 \(d_{1}\) 诱导的映射。因此 \(H_{1}(x,M)\) 的消失等价于 \((\alpha_{M}^{x})_{1}\) 的单射性，从而得到蕴含关系 (iv) \(\Rightarrow\) (iii)。

为证明(iii)蕴含(iv)，我们将使用以下引理：

引理 4. --- 设 \(\mathbf{A}_{0}\) 为环 \(\mathbf{Z}[\mathrm{T}_{1},\ldots,\mathrm{T}_{n}]\) ，并设 \(u:\mathbf{A}_{0}[\mathrm{X}_{1},\ldots,\mathrm{X}_{n}] \to \mathbf{A}_{0}[\mathrm{U}]\) 为 \(\mathbf{A}_{0}\) -代数的、使得 \(u(\mathbf{X}_{i}) = \mathbf{T}_{i}\) U 的同态。\(u\) 的核是理想 \(\mathfrak{d}_{0}\) ，它由 \(\mathbf{A}_{0}[\mathbf{X}_{1},\ldots,\mathbf{X}_{n}]\) 中形如（ \(\mathrm{T}_{i}\mathbf{X}_{j}-\mathrm{T}_{j}\mathbf{X}_{i}\) ）的元素（ \(1 \leqslant i < j \leqslant n\) ）生成。若 t 是 \(\mathbf{A}_{0}\) 中由 ( \(\mathrm{T}_{1},\ldots,\mathrm{T}_{n}\) ) 生成的理想，则 \(u\) 诱导出同构

\[
\overline {{u}}: \mathrm{A} _ {0} [ \mathrm{X} _ {1}, \dots , \mathrm{X} _ {n} ] / \mathfrak {d} _ {0} \rightarrow \bigoplus_ {r \geqslant 0} t ^ {r}.
\]

显然只需证明第一个断言。对每个自然整数序列 \(\alpha = (\alpha_{1},\ldots,\alpha_{n})\) 及每个整数 \(k\in[0,n]\) ，记 \(\mathbf{P}_{\alpha,k}\) 为单项式

\[
\mathrm{T} _ {1} ^ {\alpha_ {1}} \dots \mathrm{T} _ {k} ^ {\alpha_ {k}} \mathrm{X} _ {k + 1} ^ {\alpha_ {k + 1}} \dots \mathrm{X} _ {n} ^ {\alpha_ {n}};
\]

令 N 为 \(A_{0}[X_{1}, \ldots, X_{n}]\) 的子 Z-模，由 \(P_{\alpha,k}\) 对 \(\alpha \in N^{n}\) 和 \(0 \leqslant k \leqslant n\) 生成。我们将证明 u 在 N 上的限制是单射，且 \(A_{0}[X_{1}, \ldots, X_{n}] = \mathfrak{d}_{0} + N\) ；由于有 \(d_{0} \subset Ker u\) ，由此即得引理。

注意 N 由集合 S 生成，S 由 \(P_{0,0} = 1\) 以及那些 \(P_{\alpha,k}\) （对它们有 \(\alpha_k \neq 0\) ）组成。为证明 u 在 N 上的限制是单射的，只需证明 S 中两个不同的元素在 u 下的像是 \(A_0[U]\) 中不同的单项式。现在有 \(u(P_{\alpha,k}) = T^\alpha U^{\sum_{i \geq k} \alpha_i}\) ，因此等式 \(u(P_{\alpha,k}) = u(P_{\alpha',k'})\) 蕴含 \(\alpha = \alpha'\) 与 \(\sum_{i \geq k} \alpha_i = \sum_{i \geq k'} \alpha_i\) 。假设 \(P_{\alpha,k}\) 和 \(P_{\alpha',k'}\) 属于 S。若 \(\alpha = 0\) ，则有 \(k = k' = 0\) ；若 \(\alpha \neq 0\) ，我们得到 \(k = k'\) ，因为 \(\alpha_k\) 与 \(\alpha_{k'}\) 非零，从而结论成立。

让我们证明每个单项式 \(\mathbf{T}^{\alpha}\mathbf{X}^{\beta}\in\mathbf{A}_{0}[\mathbf{X}_{1},..., \mathbf{X}_{n}]\) 模 \(\mathfrak{d}_{0}\) 同余于某个 \(\mathbf{P}_{\eta,k}\) 。对每个序列 \(\lambda\in\mathbf{N}^{n}\) ，我们用 \(i(\lambda)\) （相应地， \(j(\lambda)\) ）表示最小的（相应地，最大的）整数 \(k\in[1,n]\) ，它使得 \(\lambda_{k}\neq0\) 。在单项式 \(\mathbf{T}^{\gamma}\mathbf{X}^{\delta}\) （它们与 \(\mathbf{T}^{\alpha}\mathbf{X}^{\beta}\) 模 \(\mathfrak{d}_{0}\) 同余）的集合中，我们选择一个使有理整数 \(j(\gamma)-i(\delta)\) 为最小的；让我们证明此时有 \(j(\gamma)-i(\delta)<0\) 。假设我们有 \(j(\gamma)\geqslant i(\delta)\) ；记 \(j=j(\gamma)\) , \(i=i(\delta)\) ，并设 \(\varepsilon=\inf(\gamma_{j},\delta_{i})\) 。由于 \((\mathbf{T}_{j}^{\varepsilon}\mathbf{X}_{i}^{\varepsilon}-\mathbf{T}_{i}^{\varepsilon}\mathbf{X}_{j}^{\varepsilon})\) 能被 \((\mathbf{T}_{j}\mathbf{X}_{i}-\mathbf{T}_{i}\mathbf{X}_{j})\) 整除，故属于 \(\mathfrak{d}_{0}\) ，我们看到 \(\mathbf{T}^{\gamma}\mathbf{X}^{\delta}\) 模 \(\mathfrak{d}_{0}\) 同余于 \(\mathbf{T}^{\gamma'}\mathbf{X}^{\delta'}\) ，其中 \(\gamma_{i}'=\gamma_{i}+\varepsilon\) , \(\gamma_{j}'=\gamma_{j}-\varepsilon\) , \(\gamma_{k}'=\gamma_{k}\) （ \(k\neq i,j\) ），且 \(\delta_{i}'=\delta_{i}-\varepsilon\) , \(\delta_{j}'=\delta_{j}+\varepsilon\) , \(\delta_{k}'=\delta_{k}\) （ \(k\neq i,j\) ）。由于 \(\gamma_{j}'\) 或 \(\delta_{i}'\) 为零，我们有 \(j(\gamma')-i(\delta')<j(\gamma)-i(\delta)\) ，这与 \(j(\gamma)-i(\delta)\) 的极小性相矛盾。

因此我们有 \(j(\gamma) < i(\delta)\) ，从而 \(\mathbf{T}^{\gamma}\mathbf{X}^{\delta} \in \mathbf{N}\) ，这就完成了引理的证明。

让我们证明 (iii) 蕴含 (iv)。考虑环 \(\mathbf{A}_0 = \mathbf{Z}[\mathrm{T}_1, \ldots, \mathrm{T}_n]\) 以及 \(\mathbf{A}_0\) 中由 \(\mathrm{T}_1, \ldots, \mathrm{T}_n\) 生成的理想 t。赋予 M 一个 \(\mathbf{A}_0\) -模结构，使得 \(\mathrm{T}_i m = x_i m\) 对 \(m \in \mathbf{M}\) , \(1 \leqslant i \leqslant n\) 成立。根据 X，第 155 页， \(\mathrm{H}_i(x, \mathbf{M})\) 与 \(\mathbf{H}_i(\mathbf{T}, \mathbf{M})\) 典范地等同。

由于序列 \(\mathbf{T}\) 对 \(\mathbf{A}_{0}\) 是正则的，由蕴含关系 (i) \(\Rightarrow\) (ii) 可知，复形 \(\mathbf{K}.\left(\mathbf{T},\mathbf{A}_{0}\right)\) 定义了 \(\mathbf{A}_{0}\) -模 \(\mathbf{A}_{0}/t\) 的一个自由分解。注意，该模与 \(\mathbf{Z}\) 等同，配备 \(\mathbf{A}_{0}\) -模结构，使得 \(\mathrm{T}_{i}\mathbf{Z}=0\) 对 \(1\leqslant i\leqslant n\) 成立。因此条件 (iii) 等价于 \(\operatorname{Tor}_{1}^{\mathbf{A}_{0}}(\mathbf{M},\mathbf{Z})=0\) 。

让我们证明 (iii) 蕴含 \(\mathrm{Tor}_{1}^{\mathbf{A}_{0}}(\mathbf{M},\mathbf{A}_{0}/t^{r})=0\) 对所有 \(r\geqslant1\) 成立。这对 \(r=1\) 的情形由前面的论述得出。在一般情形，考虑正合列

\[
0 \rightarrow t ^ {r} / t ^ {r + 1} \rightarrow A _ {0} / t ^ {r + 1} \rightarrow A _ {0} / t ^ {r} \rightarrow 0.\tag{16}
\]

该 \(A_{0}\) -模 \(t^{r}/t^{r+1}\) 同构于 Z 的有限份拷贝的直积；因此我们有 \(\operatorname{Tor}_{1}^{\mathbf{A}_{0}}(\mathbf{M}, t^{r}/t^{r+1}) = 0\) 。由此，通过对 r 归纳，我们推出

\[
\operatorname{Tor} _ {1} ^ {\mathrm{A} _ {0}} \left(\mathrm{M}, \mathrm{A} _ {0} / t ^ {r}\right) = 0 \quad \text { for   all } r.
\]

因此，正合列 (16) 通过与 M 作张量积，给出一个正合列

\[
0 \rightarrow (t ^ {r} / t ^ {r + 1}) \otimes_ {\mathrm{A} _ {0}} \mathrm{M} \rightarrow \mathrm{M} / x ^ {r + 1} \mathrm{M} \rightarrow \mathrm{M} / x ^ {r} \mathrm{M} \rightarrow 0\tag{17}
\]

从而得到 \((\mathfrak{t}^r /\mathfrak{t}^{r + 1})\otimes_{\mathbf{A}_0}\mathbf{M}\) 到 \(\mathfrak{x}^r\mathbf{M} / \mathfrak{x}^{r + 1}\mathbf{M}\) 上的一个同构。考虑正合列 \(0\to \mathfrak{t}^{r + 1}\to \mathfrak{t}^r\to \mathfrak{t}^r /\mathfrak{t}^{r + 1}\to 0\) ，然后对 \(r\) 归纳可以看出，映射 \(m_r:\mathfrak{t}^r\otimes_{\mathbf{A}_0}\mathbf{M}\rightarrow \mathfrak{x}^r\mathbf{M}\) （它由 \(\mathbf{A}_0\) 在 \(\mathbf{M}\) 上的作用诱导）是一个同构。

为证明 (iv)，只需观察到图表

\[
\begin{array}{c} \left(\mathrm{A} _ {0} \left[ \mathrm{X} _ {1}, \dots , \mathrm{X} _ {n} \right] / \mathfrak {d} _ {0}\right) \otimes_ {\mathrm{A} _ {0}} \mathrm{M} \xrightarrow {\bar {u} \otimes 1 _ {\mathrm{M}}} \bigoplus_ {r \geqslant 0} \left(\mathrm{t} ^ {r} \otimes_ {\mathrm{A} _ {0}} \mathrm{M}\right) \\ \Bigg \downarrow^ {e} \\ ^ {\prime} (\mathrm{A} [ \mathrm{X} _ {1}, \dots , \mathrm{X} _ {n} ] / \mathfrak {d}) \otimes_ {\mathrm{A}} \mathrm{M} \xrightarrow {\alpha_ {\mathrm{M}} ^ {*}} \bigoplus_ {r \geqslant 0} x ^ {r} \mathrm{M} \end{array}
\]

其中 e 是纯量扩张的典范同构（II，第 85 页，命题 6），该图表是交换的，然后应用引理 4 即可。

